\section{引言}
格点 QCD 模拟所预言的、与粒子和强子物理唯象学相关的可观测量正日益增多。
这些结果通常从大欧氏时间间隔处 $n$ 点函数的期望值中提取。
由于这些函数随时间衰减，当时间间隔取得很大时，
统计噪声与信号之比呈指数式增大（零动量赝标介子是一个显著的例外）。
幸运的是，在构造用于产生介子态和重子态的内插算符时存在一定的自由度。
采用与基态波函数空间结构相似的内插算符，
可以在信噪比仍然很大的时间间隔处提取渐近结果。

如今许多应用需要携带动量的强子。
例如，将 $B\rightarrow\pi\ell\bar{\nu}_{\ell}$ 或
$\Lambda_b\rightarrow p\ell\bar{\nu}_{\ell}$
半轻子衰变形状因子的计算~\cite{Detmold:2015aaa}
推向小虚度区域，分别要求空间动量达到
$B$ 介子与 $\pi$ 介子之间或 $\Lambda_b$ 重子与质子之间质量差的大小。
另一类重要应用是部分子分布函数（PDF）及其推广，
特别是横动量依赖部分子分布函数（TMD），以及作为整体的 Wigner 分布。
这些量同样是从 $ \bra{H(\vect{p}')} \oper \ket{H(\vect{p})}$ 型的矩阵元中提取的，其中 $H(\vect{p})$
是动量为 $\vect{p}$ 的强子态。
算符 $\oper$ 在欧氏格点上不能有闵可夫斯基时间方向的延展。因此，为了对一个内禀非定域的
算符外推到光前运动学区域，
必须在强子携带高空间动量 $\vect{p}$、$\vect{p}'$ 的参考系中进行计算。
这一事实早在格点 TMD 计算的背景下就已为人所知~\cite{Hagler:2009mb,Musch:2010ka,Musch:2011er}，
并且在最近关于 $\pi$ 介子中这些分布的工作中得到了非常清晰的说明~\cite{Engelhardt:2015xja}。
出于同样的原因，在一个新提出的、以受控方式将准部分子分布（quasi PDF）
与光前分布联系起来的方案中，
格点上快速运动的强子也是十分理想的~\cite{Ji:2013dva,Ji:2015jwa}。
沿此方向的首批格点
计算已经开始~\cite{Lin:2014zya,Alexandrou:2015rja}。
早先提出的在位置空间中计算准分布振幅的建议~\cite{Aglietti:1998ur,Abada:2001if,Braun:2007wv}
同样需要
高动量的 $\pi$ 介子或核子。遗憾的是，
迄今为止尚无令人满意的技术能在强子携带高动量时
抑制激发态贡献。

更具体地说，两点函数由下式给出：
\begin{equation}
\label{eq:twopoint}
C_H(t)=\langle H(t)H^{\dagger}(0)\rangle
=\sum_{j>0}\frac{|\langle j|\hat{H}^{\dagger}|0\rangle|^2}{2E_{H,j}}e^{-E_{H,j}t}\,,
\end{equation}
其中 $E_{H,j}$ 表示由内插算符
$H^{\dagger}$ 所产生的态塔中的第 $j$ 个能级。显然，第 $j$ 个激发态的贡献
不仅相对于基态受到 $\exp[-(E_{H,j}-E_{H,1})t]$ 的压制，
还受到比值
$|\langle j|\hat{H}^{\dagger}|0\rangle|^2/|\langle 1|\hat{H}^{\dagger}|0\rangle|^2$
的压制：增大基态重叠因子
$|\langle 1|\hat{H}^{\dagger}|0\rangle|$
相对于 $j>1$ 时 $|\langle j|\hat{H}^{\dagger}|0\rangle|$ 的比值，
会导致对激发态的额外压制。

通过采用延展的内插算符来减小激发态重叠，最早是在
胶球谱的计算中进行的。在此情形下，可以对相应内插算符中的规范链接进行迭代
``APE 涂抹''~\cite{Falcioni:1984ei}、
``fuzzing（绒化）''~\cite{Teper:1987wt} 或 ``HYP 涂抹''~\cite{Hasenfratz:2001hp}，
以更好地逼近（光滑的）基态波函数，
另见文献~\cite{Gupta:1990ek}。
这种规范链接涂抹随后被推广到内插算符中夸克场的迭代涂抹，
特别是规范协变的 Wuppertal（即高斯）涂抹~\cite{Gusken:1989ad,Gusken:1989qx,Alexandrou:1990dq}、
类氢涂抹~\cite{Alexandrou:1990dq}
以及 Jacobi 涂抹~\cite{Collins:1992rp,Allton:1993wc}。
此外，文献~\cite{Gusken:1989qx,Alexandrou:1990dq,Bali:2005fu}
在夸克涂抹中对空间规范输运子采用了 APE 涂抹，而文献~\cite{Bali:2005fu}
则利用了不同 Wuppertal 涂抹次数的线性组合。

由于高斯型涂抹函数即使在内插算符中加入导数之后也可能不适合产生例如
$p$ 波态，迭代涂抹后来在文献~\cite{Lacock:1994qx} 中与位移夸克源（fuzzing）
相结合，文献~\cite{Boyle:1999gx} 提出了其推广。
最后，文献~\cite{vonHippel:2013yfa} 发明了``自由形式涂抹''，
以规范协变的方式将高斯涂抹与任意函数折叠。
在规范协变迭代涂抹函数之前及其同时，也使用过固定规范的源：
用于零动量的壁源~\cite{Bacilieri:1988fh} 和
非零动量的壁源~\cite{Detmold:2012wc}、盒状源~\cite{Allton:1990qg}、
高斯型``壳层源''~\cite{DeGrand:1990dz}
以及带节点的源~\cite{DeGrand:1991ng}。
这些固定规范的方法与自由形式涂抹共同的缺点是：
涂抹汇点（sink）时需要对所有夸克位置逐一求和，
代价高昂到不可行。然而，让源与汇采用相同的内插算符是非常理想的，
因为只有这样才能保证谱分解
式~\eqref{eq:twopoint} 中系数的正性，
从而保证两点函数的凸性。为完整起见，
我们还提及文献~\cite{Peardon:2009gh} 的 ``distillation''（蒸馏，即
Laplacian-Heaviside）方法，
因为它与规范协变涂抹密切相关。

大动量增加了态的能量，使两点函数和三点函数衰减更快，
因而信噪比更差。此外，正如我们将看到的，
使用传统夸克涂抹方法时激发态压制的效果会大大降低。
已有一些尝试~\cite{Roberts:2012tp,DellaMorte:2012xc}
在 Wuppertal 涂抹~\cite{Gusken:1989ad,Gusken:1989qx} 中引入各向异性，
目的是让内插波函数沿空间
动量 $\vect{p}$ 的方向按 boost 因子 $1/\gamma=m/E(\vect{p})$
进行洛伦兹收缩。然而，
这并未带来人们所期待的基态增强。这里我们将论证并证明：要在高动量下获得满意的
结果，需要在夸克涂抹函数中加入额外的相位因子。

本文结构如下。
第~\ref{sec:basic} 节讨论我们所引入的新一类涂抹函数背后的基本思想。
第~\ref{sec:wuppertal} 节更加具体，以 Wuppertal 涂抹作为一般性例子加以修改并提出进一步改进。
第~\ref{sec:simulate} 节讨论我们的模拟参数以及
对核子和 $\pi$ 介子能量的预期。铺垫完成之后，
第~\ref{sec:results} 节通过现实的数值研究考察该方法的可行性，
优化涂抹参数并对新旧方法进行比较。最后，
我们研究 $\pi$ 介子和核子的色散关系，然后作结论。
